\section{Forgotten Tradition}

There seems to be much misunderstanding of Evola's moral position. As he claims, he is drawing on

\begin{quotex}
doctrines which are in truth of a fundamental value, currently almost forgotten, that placed evil in matter.
\end{quotex}


Apparently contemporary readers really have forgotten them, even those who falsely believe themselves to be relying on Evola in their Traditionalism and paganism. Perhaps such men need to rethink their relationship to Evola. As an example of such an attitude, we need do no more than quote the first line of Porphyry's biography of Plotinus:

\begin{quotex}
Plotinus, the philosopher our contemporary, seemed ashamed of being in the body.
\end{quotex}


In sharp contrast to that, we need look no further than the recent Hollywood film, \emph{The Little Fockers} to see the modern mind at work. In it, the European family patriarch is portrayed as humourless and up-tight. Finally, he is defeated by his son-in-law, after which he is informed that he is really a Jew. The benefit of that, as he is told, is that farting, belching, sexing and all such bodily functions --- things that used to embarrass the patriarch --- are now ``good''.

Contrast that to the Traditional, aristocratic attitude of a Plotinus. Make up your mind once and for all where you stand: with Tradition or with the modern world.

Henry Chadwick writes in \textit{Interpreting Late Antiquity}\footnote{\url{https://www.amazon.com/Interpreting-Late-Antiquity-Essays-Postclassical/dp/0674005988/ref=sr\_1\_1?ie=UTF8&qid=1300597497&sr=8-1}}: 

\begin{quotationx}
Porphyry, biographer and editor of Plotinus, began the Life of his master with the famous sentence that Plotinus always seemed ashamed of being in the body. His concern was with the soul which, true to Plato, he saw as being midway between the inferior flesh and the superior incorporeal mind that is, spirit . \emph{Moral choices are therefore decisions whether to follow the higher reason or lower bodily appetite.} The self (for which Plotinus uses the word \emph{autos}) therefore has levels and power of movement. But the true self is divine and the body no more than its temporary instrument (IV 7.1.20ff.).

Naturally the union of soul with body makes for interaction between them. When ashamed we blush, when terrified we go white with fear. Conversely there are pressures which the body can put upon the soul, drawing it downward toward material satisfactions. \emph{But the soul's true home is in that higher realm which is not governed by the determinism of fate.} Self-knowledge is synonymous with the soul's being identical in being (\emph{homoousios}) with mind (IV 4.28.56; IV 7.10.19). Thereby \emph{salvation is the divinizing of the soul}, a mystical union comparable to the merging of two torches ( 6.9.18ff.) ``in the measure possible to the human soul'' (12-9-9.45ff.). The route by which one returns to true being is self-knowledge (VI 5.7), and is a restoration of the unity from which the soul has fallen into multiplicity and has been torn apart in a ``scattering'' (VI 6.I.5).
\emph{As the soul moves toward the good, it recovers freedom, which is a liberation from the constraints of the body} (VI 8.6-7). Plotinus noted explicitly that in his doctrine of the soul's retaining its divine nature ``undescended'' he was departing from the normal view of the Platonic schools (IV 8.8.I). Origen (Prin. 3.4) shows that the idea was not new with Plotinus.\footnote{All emphases are mine.}
\end{quotationx}



\flrightit{Posted on 2011-03-20 by Cologero}

\begin{center}* * *\end{center}

\begin{footnotesize}
\begin{sffamily}
\texttt{EXIT on 2011-03-20 at 08:12 said:}

I don't care who wrote what. I don't base principles on someone's emotions, especially shame of the body. Do you actually expect me to accept Plotinus' alleged shame as proof that matter is evil? But of course you say you do not wish to convert anyone yet you urge us to rethink our positions? I don't rely on Evola for my tradition, I rely on my Self and my Body.

\hfill

\texttt{EXIT on 2011-03-20 at 08:22 said:}

Is it not written in the New Testament that the body is a temple?

\hfill

\texttt{EXIT on 2011-03-20 at 08:34 said:}

1 Cor. 6:18-20, ``Every sin that a man doeth is without the body; your body is the temple of the Holy Ghost which is in you... therefore glorify God in your body and spirit, which are God's.''

\hfill

\texttt{James O'Meara on 2011-03-20 at 12:02 said:}

I take your basic point, but based on a nodding acquaintance with both Evola and Neo-Platonism, I don't think Evola would have considered Plotinus' shame even assuming it to be authentic, rather than another of Porphyry's mythologizing touches, like the snake that slithers from under Plotinus' deathbed, or the nine-fold arrangement of his Treatises to be a positive aspect of the Aryan self-image. The classical Greeks and Romans, on whom Evola based his idea of the ideal Roman type see Revolt and even the later Men/Ruins , were not `promiscuous' like the ``mandolin strumming, macaroni scarfing'' Mediterranean types, but hardly `ashamed' of their bodies.

I, most scholars, and I think Evola, would take that notion as a symptom of Late antiquity, either under the influence of Christianity the `primitive' type, not the Roman, in Evola's terms or else of a weakening of the Roman spirit itself, which is what made the influencing possible.

Like your previous remarks on sola sciptura resulting in the demotion of Classical myth, religion and philosophy from `revealed' status and thus cutting Christianity off from its roots in Tradition, some, mostly the same kind of Christians, would agree that Hellenism was different from Classical, but see this as a plus for both pagans and Christianity the `moral advance over the gross pagans' idea , while Evola would deplore it. But either way, it IS different.

\hfill

\texttt{Cologero on 2011-03-20 at 13:01 said:}

Good points, James, yet we must be very careful about how we read deep thinkers. To be ashamed of one's body is no sign of spirituality, but apparently Plotinus was ashamed of ``being'' in a body, not precisely the same notion. My intention was not to promote such shame, but rather to point out that Evola's contention --- that ``necessity'' is inherently an ``evil'' --- is not unique to him, but has deeper roots in tradition. Instead of Porphyry, focus instead on the commentary by Chadwick that I quoted. So Plotinus' shame was not the body \emph{per se}, but rather the necessity that limited his freedom and power.

Now Evola does indeed address the idea of ``aristocratic asceticism'' in Revolt (ch 16), where he uses the towering figures of Plotinus, Buddha, and Meister Eckhart to support his point. This asceticism involves a detachment from all things (yes, the ``body'' is a thing). Evola writes: ``An action dictated by desire ... must not be undertaken.'' Of course, the ignorant and prudish assume this means that the specific desire or the object of desire is evil. Not at all, that misses the point entirely. The issue at stake is the feeling of ``necessity'' that \emph{is} desire; that is what must be overcome.

Looking at chapter 1 of Revolt, Evola writes: ``Hellenism saw nature as the embodiment of eternal ``privation'' yes, that is the correct translation, relating it back to what we have previously written I'll expect the reader to turn to the rest of that paragraph. Note, too, that Evola relates this also back to the Scholastics, something we tried to show in the chart.

At the very beginning of Revolt, Evola makes clear that he is speaking from ``knowledge'' not mere theory. We can speculate all we want about the self-image of the Aryan, his qualities, attitudes, and so on. Or else, we can actually ``be'' an Aryan (twice-born) in Spirit and know first hand.

\hfill

\texttt{James O'Meara on 2011-03-20 at 15:07 said:}

All true enough, but I only wanted to emphasize that we can't really rely on Porphyry; he's not writing what we would call a `biography' and in fact nothing like that even existed at the time. When Porphyry says ``Plotinus seemed ashamed'' we are at least 2 steps away from Plotinus himself; what Plotinus thought, how he seemed to Porphyry, how Porphyry fits this into his ``Life of the Sage'' template. Now, when I compare that language to Plotinus's own, it seems suspicious; its sounds more like the Christians he defines himself against in Against the Gnostics the only Christian writers he seems to know, interestingly as ungrateful wretches daring to despise the world which, defective as it is, his point and yours is none the less the best possible work of the World Soul. Of course, what you, Evola and Chadwick present is a different matter, since it's based on Plotinus himself.

As for theory and `knowledge,' I'm not sure where it applies here, as the Aryan self-image I reference is also from Evola, from the same book.

Also, I forgot to mention the Fockers reference is spot on; on my blog I have frequently discussed the anti-European `cockroach literature' promoted by the media and academy, based on Kevin MacDonald's studies of the ``Culture of Critique'' and Cuddihy's ``Ordeal of Civility.''

\hfill

\texttt{Francis on 2011-03-20 at 21:50 said:}

Although EXIT is hardly inviting conversation, perhaps this is due to genuine concerns of which I am unaware. I am a young'in and fairly new to here so forgive my lack of decent input.

The `body' certainly fascinates an awkward position in the cosmos and so begs a fair amount of attention. That with it time and space occupy from nil; that eternal and terminal somehow meet with the Logos; that the reason for the hypostatic union is the failure of the being with a body... all of this asks the place of the body as it is now, before death.

While the body is a temple, and EXIT does not wish to deny it, that temple will break down and die, and all will have to deny it at death, less they choose to stay with it in death.

1Cor6:18-20 must be reconciled with this and Rom6:6, Rom7:24, Rom8:10, Rom8:13... Not to mention, 1Cor6:18-20 does not necessarily carry with it the meaning you are applying. I was going to grant you that it might be the literal meaning... but even the literal meaning is about the `body' of Christians who are to be ``members of Christ'' and not in cahoots with other `bodies' (ie. there is discord/division/disagreement in the community). In the next section, Paul then addresses issues related to the body. And there, Paul is decidedly negative on the sexual deeds of the body... only good for those who cannot control their passions.

Does this help?

I do not have the time now, but I think Chadwick helpfully refers to Origen, who is also helpful on this point. His `speculation' (spiritual insight?) in ``first principles'' on the `body,' places the body in a simultaneously significant and insignificant position.

\hfill

\texttt{HOO on 2011-03-20 at 23:25 said:}

Very good, James O'Meara, thank you.

EXIT, amusing as ever. But why don't you try to be more positive? I see you cited Corinthians, ´your body is the temple´, there. That's very good! Now if you would take more time to elaborate on your thoughts, that would be even better. You'll notice I sometime include multiple citations like that from various works, in my posts. Why don't you be a buddy and do more of that rather than mostly ´criticizing´.

I think I agree with your stance on the body. In any case I would think everybody here recognizes that there are at least two ways of dealing with it. One way is ´life-denying´ and ´neglects´ the body for trans-worldly states, the other way is ´(super)life-affirming´ and ´celebrates´ the body for this-worldly and trans-worldly states.

``We are mind and body: if mind and body (inasmuch as they belong to the world of maya) are false, how can one hope to achieve through them that which is true?'' quoted in ``The Yoga of Power'' p. 19 

´The analysis of the last age, the ``dark age'' or Kali Yuga, brings to light two essential features. The first is that mankind living in this age is strictly connected to the body and cannot prescind from it; therefore, the only way open is not that of pure detachment (as in early Buddhism and in the many varieties of yoga) but rather that of knowledge, awakening, and mastery over secret energies trapped in the body.´ Evola, ``The Yoga of Power'' p. 3 ´In Tantrism the passage from the ideal of ``liberation'' to that of ``freedom'' marks an essential change in the ideals and ethics of Hinduism. It is true that even previously the ideal of the jivanmukta had been known. The word means ``one who is freed,'' that is, the one who has achieved the unconditioned, the sahaja, while alive, in his own body.´ Evola, ``The Yoga of Power'' p. 4 The prechristian world knew this in different traditions. It was known by some in the Christian era too. And it is known by some in the ´postchristian´ era.

Basic stuff, guys.

\hfill

\texttt{Cologero on 2011-03-21 at 01:26 said:}

This discussion has gotten out of hand, as the original point has been lost. Evola wrote:

``evil in matter, meaning by this that residue of necessity that resists the idea --- that is, freedom --- and limits it''

Unless one understands this definition of matter, the point is lost. Nowhere does Evola state the body in itself is evil.

Insofar as the body is matter, it sets a limit to freedom, that is, it is privation. But, insofar as the sequence of domination Spirit-> Soul-> Body holds, the body is not ``evil''.

I then tied this into a larger framework. First the Scholastics, who Evola references in Revolt. If the fall is the weakening and bondage of the will, then it would seem that Power and Freedom are Good. I used Plotinus to illustrate, and was then forced to show that Plotinus is not an aberration of late antiquity, since Evola is in fundamental with him.

It would be of value to start again at the beginning and follow the argument, point by point.

\hfill

\texttt{Zmaj on 2011-03-21 at 06:54 said:}

Cologero has explained the case well, I see no issue. To dismiss a few possible confusions that could be stoking disputes:

-Use of the term `evil' in these articles should not evoke any moralistic connotations. In the same way that `evil' implies a subversion or counterfeit in theology, we might employ the term `negate' to describe the same tendency in metaphysic. This clarified, I'll retain the term `evil' hereinafter.

-Matter, as such, is evil, especially insofar as it is understood according to the modern conception (which conception negates everything other than matter). To the extent that any being on the temporal plane has simultaneous access to higher perception, the evil of matter is mitigated because said being can begin to perceive the general essence underlying the manifestation of matter. Indeed, the modern idea of matter is revealed as un erreur majeur to anyone who breaks the barrier.

-This brings us to comments like ``The body is a temple'', and discussions about yoga. If the body is understood as a point of departure in supratemporal realization, it is no longer just plain material ideas that we are dealing with. So, it can be a temple that grounds the center of consciousness until sufficient energies begin to displace that center toward its transcendent origin. Thereafter, the idea is to become less dependent on the ephemeral, body included -- one awakens to oneself in the body; this does not make the body the origin.

-The body as matter is also evil in the sense that it can condemn your soul by trapping it and limiting realization. Recall certain physical conditions that are prerequisites of some yogic praxis.

\hfill

\end{sffamily}
\end{footnotesize}

